% This contents of this file will be inserted into the _Solutions version of the
% output tex document.  Here's an example:

% If assignment with subquestion (1.a) requires a written response, you will
% find the following flag within this document: <SCPD_SUBMISSION_TAG>_1a
% In this example, you would insert the LaTeX for your solution to (1.a) between
% the <SCPD_SUBMISSION_TAG>_1a flags.  If you also constrain your answer between the
% START_CODE_HERE and END_CODE_HERE flags, your LaTeX will be styled as a
% solution within the final document.

% Please do not use the '<SCPD_SUBMISSION_TAG>' character anywhere within your code.  As expected,
% that will confuse the regular expressions we use to identify your solution.

\def\assignmentnum{1 }
\def\assignmentname{Sentiment Analysis}
\def\assignmenttitle{XCS221 Assignment \assignmentnum --- \assignmentname}
\input{macros}
\begin{document}
\pagestyle{myheadings} \markboth{}{\assignmenttitle}

% <SCPD_SUBMISSION_TAG>_entire_submission

This handout includes space for every question that requires a written response.
Please feel free to use it to handwrite your solutions (legibly, please).  If
you choose to typeset your solutions, the |README.md| for this assignment includes
instructions to regenerate this handout with your typeset \LaTeX{} solutions.
\ruleskip

\LARGE
1.e


\normalsize

% <SCPD_SUBMISSION_TAG>_1e
\begin{answer}
  % ### START CODE HERE ###
  1- The validation error is around 0.27 for word features and n-gram when n is around 7. This is due to the fact that with n-grams of size 7, the model can capture more context and semantic meaning from the text, which helps in better understanding. Smaller n-grams may miss important phrases, while larger n-grams may lead to sparsity issues. Therefore, n-grams of around 7 strike a good balance between capturing context and maintaining sufficient data for training.
  
  2- When a user writes a review with informal language like "Thiiis movie was suuuuuuuper scaaryyy!!!!" Because the review contains misspellings and stretched words (“Thiiis”, “suuuuuuuper”, “scaaryyy”), word features treat each as a completely new and unseen token, while character n-grams still capture the emotional patterns inside the words (e.g., “thi”, “hii”, “iii”, "iis"), and will still identify the word as positive or negative based on the presence of these patterns. Therefore, character n-grams are more robust to such variations in spelling and can better capture the sentiment of the review.
  % ### END CODE HERE ###
\end{answer}
% <SCPD_SUBMISSION_TAG>_1e
\clearpage

\LARGE
2.a
\normalsize

% <SCPD_SUBMISSION_TAG>_2a
\begin{answer}
% ### START CODE HERE ###

We run K=2 means on the four points:
\[
x_1=(10,0),\;
x_2=(30,0),\;
x_3=(10,20),\;
x_4=(20,20).
\]

\subsection*{Case (i): Initialization $\mu_1^{(0)} = (20,30)$, $\mu_2^{(0)} = (20,-10)$}

\textbf{Step 1 (Set assignments).}

For each point index $i = 1,\ldots,4$ we assign
\[
z_i^{(1)}
=
\arg\min_{k=1,2}
\,\bigl\|\, x_i - \mu_k^{(0)} \,\bigr\|^2.
\]

Example (computing squared distances):
\[
\|x_1-\mu_1^{(0)}\|^2 = 1000,
\qquad
\|x_1-\mu_2^{(0)}\|^2 = 200.
\]
Thus
\[
z_1^{(1)} = 2.
\]

Similarly:
\[
z_2^{(1)} = 2,\quad
z_3^{(1)} = 1,\quad
z_4^{(1)} = 1.
\]

\textbf{Step 2 (Update of centroids).}

Cluster 1 points: $x_3, x_4$  
Cluster 2 points: $x_1, x_2$

Cluster sizes:
\[
m_1 = 2,\qquad m_2 = 2.
\]

Updated centroids:
\[
\mu_1^{(1)}
=
\frac{1}{2}\bigl((10,20)+(20,20)\bigr)
=
(15,20),
\]
\[
\mu_2^{(1)}
=
\frac{1}{2}\bigl((10,0)+(30,0)\bigr)
=
(20,0).
\]

\textbf{Second iteration.}

\[
z_1^{(2)} = 2,\quad
z_2^{(2)} = 2,\quad
z_3^{(2)} = 1,\quad
z_4^{(2)} = 1.
\]

Final centroids:
\[
\mu_1^{(*)} = (15,20),\qquad
\mu_2^{(*)} = (20,0).
\]


\subsection*{Case (ii): Initialization $\mu_1^{(0)} = (0,10)$, $\mu_2^{(0)} = (30,20)$}

\textbf{Step 1 (Assignments).}

\[
z_1^{(1)} = 1,\;
z_2^{(1)} = 2,\;
z_3^{(1)} = 1,\;
z_4^{(1)} = 2.
\]

\textbf{Step 2 (Update).}

\[
m_1 = 2,\qquad m_2 = 2.
\]

\[
\mu_1^{(1)}=\frac{1}{2}\bigl((10,0)+(10,20)\bigr)=(10,10),
\qquad
\mu_2^{(1)}=\frac{1}{2}\bigl((30,0)+(20,20)\bigr)=(25,10).
\]

\textbf{Final iteration.}

\[
z_1=1,\;
z_2=2,\;
z_3=1,\;
z_4=2.
\]

\[
\mu_1^{(*)}=(10,10),
\qquad
\mu_2^{(*)}=(25,10).
\]

% ### END CODE HERE ###

\end{answer}
% <SCPD_SUBMISSION_TAG>_2a
\clearpage

% <SCPD_SUBMISSION_TAG>_entire_submission
\end{document}